\documentclass[12pt,a4paper,]{book}
\def\ifdoblecara{} %% set to true
\def\ifprincipal{} %% set to true
\let\ifprincipal\undefined %% set to false
\def\ifcitapandoc{} %% set to true
\let\ifcitapandoc\undefined %% set to false
\usepackage{lmodern}
% sin fontmathfamily
\usepackage{amssymb,amsmath}
\usepackage{ifxetex,ifluatex}
%\usepackage{fixltx2e} % provides \textsubscript %PLLC
\ifnum 0\ifxetex 1\fi\ifluatex 1\fi=0 % if pdftex
  \usepackage[T1]{fontenc}
  \usepackage[utf8]{inputenc}
\else % if luatex or xelatex
  \ifxetex
    \usepackage{mathspec}
  \else
    \usepackage{fontspec}
  \fi
  \defaultfontfeatures{Ligatures=TeX,Scale=MatchLowercase}
\fi
% use upquote if available, for straight quotes in verbatim environments
\IfFileExists{upquote.sty}{\usepackage{upquote}}{}
% use microtype if available
\IfFileExists{microtype.sty}{%
\usepackage{microtype}
\UseMicrotypeSet[protrusion]{basicmath} % disable protrusion for tt fonts
}{}
\usepackage[margin = 2.5cm]{geometry}
\usepackage{hyperref}
\hypersetup{unicode=true,
            pdfauthor={Nombre Completo Autor},
              pdfborder={0 0 0},
              breaklinks=true}
\urlstyle{same}  % don't use monospace font for urls
%
\usepackage[usenames,dvipsnames]{xcolor}  %new PLLC
\IfFileExists{parskip.sty}{%
\usepackage{parskip}
}{% else
\setlength{\parindent}{0pt}
\setlength{\parskip}{6pt plus 2pt minus 1pt}
}
\setlength{\emergencystretch}{3em}  % prevent overfull lines
\providecommand{\tightlist}{%
  \setlength{\itemsep}{0pt}\setlength{\parskip}{0pt}}
\setcounter{secnumdepth}{5}
% Redefines (sub)paragraphs to behave more like sections
\ifx\paragraph\undefined\else
\let\oldparagraph\paragraph
\renewcommand{\paragraph}[1]{\oldparagraph{#1}\mbox{}}
\fi
\ifx\subparagraph\undefined\else
\let\oldsubparagraph\subparagraph
\renewcommand{\subparagraph}[1]{\oldsubparagraph{#1}\mbox{}}
\fi

%%% Use protect on footnotes to avoid problems with footnotes in titles
\let\rmarkdownfootnote\footnote%
\def\footnote{\protect\rmarkdownfootnote}


  \title{}
    \author{Nombre Completo Autor}
      \date{18/11/2021}


%%%%%%% inicio: latex_preambulo.tex PLLC


%% UTILIZA CODIFICACIÓN UTF-8
%% MODIFICARLO CONVENIENTEMENTE PARA USARLO CON OTRAS CODIFICACIONES


%\usepackage[spanish,es-nodecimaldot,es-noshorthands]{babel}
\usepackage[spanish,es-nodecimaldot,es-noshorthands,es-tabla]{babel}
% Ver: es-tabla (en: https://osl.ugr.es/CTAN/macros/latex/contrib/babel-contrib/spanish/spanish.pdf)
% es-tabla (en: https://tex.stackexchange.com/questions/80443/change-the-word-table-in-table-captions)
\usepackage[spanish, plain, datebegin,sortcompress,nocomment,
noabstract]{flexbib}
 
\usepackage{float}
\usepackage{placeins}
\usepackage{fancyhdr}
% Solucion: ! LaTeX Error: Command \counterwithout already defined.
% https://tex.stackexchange.com/questions/425600/latex-error-command-counterwithout-already-defined
\let\counterwithout\relax
\let\counterwithin\relax
\usepackage{chngcntr}
%\usepackage{microtype}  %antes en template PLLC
\usepackage[utf8]{inputenc}
\usepackage[T1]{fontenc} % Usa codificación 8-bit que tiene 256 glyphs

%\usepackage[dvipsnames]{xcolor}
%\usepackage[usenames,dvipsnames]{xcolor}  %new
\usepackage{pdfpages}
%\usepackage{natbib}




% Para portada: latex_paginatitulo_mod_ST02.tex (inicio)
\usepackage{tikz}
\usepackage{epigraph}
\input{portadas/latex_paginatitulo_mod_ST02_add.sty}
% Para portada: latex_paginatitulo_mod_ST02.tex (fin)

% Para portada: latex_paginatitulo_mod_OV01.tex (inicio)
\usepackage{cpimod}
% Para portada: latex_paginatitulo_mod_OV01.tex (fin)

% Para portada: latex_paginatitulo_mod_OV03.tex (inicio)
\usepackage{KTHEEtitlepage}
% Para portada: latex_paginatitulo_mod_OV03.tex (fin)

\renewcommand{\contentsname}{Índice}
\renewcommand{\listfigurename}{Índice de figuras}
\renewcommand{\listtablename}{Índice de tablas}
\newcommand{\bcols}{}
\newcommand{\ecols}{}
\newcommand{\bcol}[1]{\begin{minipage}{#1\linewidth}}
\newcommand{\ecol}{\end{minipage}}
\newcommand{\balertblock}[1]{\begin{alertblock}{#1}}
\newcommand{\ealertblock}{\end{alertblock}}
\newcommand{\bitemize}{\begin{itemize}}
\newcommand{\eitemize}{\end{itemize}}
\newcommand{\benumerate}{\begin{enumerate}}
\newcommand{\eenumerate}{\end{enumerate}}
\newcommand{\saltopagina}{\newpage}
\newcommand{\bcenter}{\begin{center}}
\newcommand{\ecenter}{\end{center}}
\newcommand{\beproof}{\begin{proof}} %new
\newcommand{\eeproof}{\end{proof}} %new
%De: https://texblog.org/2007/11/07/headerfooter-in-latex-with-fancyhdr/
% \fancyhead
% E: Even page
% O: Odd page
% L: Left field
% C: Center field
% R: Right field
% H: Header
% F: Footer
%\fancyhead[CO,CE]{Resultados}

%OPCION 1
% \fancyhead[LE,RO]{\slshape \rightmark}
% \fancyhead[LO,RE]{\slshape \leftmark}
% \fancyfoot[C]{\thepage}
% \renewcommand{\headrulewidth}{0.4pt}
% \renewcommand{\footrulewidth}{0pt}

%OPCION 2
% \fancyhead[LE,RO]{\slshape \rightmark}
% \fancyfoot[LO,RE]{\slshape \leftmark}
% \fancyfoot[LE,RO]{\thepage}
% \renewcommand{\headrulewidth}{0.4pt}
% \renewcommand{\footrulewidth}{0.4pt}
%%%%%%%%%%
\usepackage{calc,amsfonts}
% Elimina la cabecera de páginas impares vacías al finalizar los capítulos
\usepackage{emptypage}
\makeatletter

%\definecolor{ocre}{RGB}{25,25,243} % Define el color azul (naranja) usado para resaltar algunas salidas
\definecolor{ocre}{RGB}{0,0,0} % Define el color a negro (aparece en los teoremas

%\usepackage{calc} 


%era if(csl-refs) con dolares
% metodobib: true


\usepackage{lipsum}

%\usepackage{tikz} % Requerido para dibujar formas personalizadas

%\usepackage{amsmath,amsthm,amssymb,amsfonts}
\usepackage{amsthm}


% Boxed/framed environments
\newtheoremstyle{ocrenumbox}% % Theorem style name
{0pt}% Space above
{0pt}% Space below
{\normalfont}% % Body font
{}% Indent amount
{\small\bf\sffamily\color{ocre}}% % Theorem head font
{\;}% Punctuation after theorem head
{0.25em}% Space after theorem head
{\small\sffamily\color{ocre}\thmname{#1}\nobreakspace\thmnumber{\@ifnotempty{#1}{}\@upn{#2}}% Theorem text (e.g. Theorem 2.1)
\thmnote{\nobreakspace\the\thm@notefont\sffamily\bfseries\color{black}---\nobreakspace#3.}} % Optional theorem note
\renewcommand{\qedsymbol}{$\blacksquare$}% Optional qed square

\newtheoremstyle{blacknumex}% Theorem style name
{5pt}% Space above
{5pt}% Space below
{\normalfont}% Body font
{} % Indent amount
{\small\bf\sffamily}% Theorem head font
{\;}% Punctuation after theorem head
{0.25em}% Space after theorem head
{\small\sffamily{\tiny\ensuremath{\blacksquare}}\nobreakspace\thmname{#1}\nobreakspace\thmnumber{\@ifnotempty{#1}{}\@upn{#2}}% Theorem text (e.g. Theorem 2.1)
\thmnote{\nobreakspace\the\thm@notefont\sffamily\bfseries---\nobreakspace#3.}}% Optional theorem note

\newtheoremstyle{blacknumbox} % Theorem style name
{0pt}% Space above
{0pt}% Space below
{\normalfont}% Body font
{}% Indent amount
{\small\bf\sffamily}% Theorem head font
{\;}% Punctuation after theorem head
{0.25em}% Space after theorem head
{\small\sffamily\thmname{#1}\nobreakspace\thmnumber{\@ifnotempty{#1}{}\@upn{#2}}% Theorem text (e.g. Theorem 2.1)
\thmnote{\nobreakspace\the\thm@notefont\sffamily\bfseries---\nobreakspace#3.}}% Optional theorem note

% Non-boxed/non-framed environments
\newtheoremstyle{ocrenum}% % Theorem style name
{5pt}% Space above
{5pt}% Space below
{\normalfont}% % Body font
{}% Indent amount
{\small\bf\sffamily\color{ocre}}% % Theorem head font
{\;}% Punctuation after theorem head
{0.25em}% Space after theorem head
{\small\sffamily\color{ocre}\thmname{#1}\nobreakspace\thmnumber{\@ifnotempty{#1}{}\@upn{#2}}% Theorem text (e.g. Theorem 2.1)
\thmnote{\nobreakspace\the\thm@notefont\sffamily\bfseries\color{black}---\nobreakspace#3.}} % Optional theorem note
\renewcommand{\qedsymbol}{$\blacksquare$}% Optional qed square
\makeatother



% Define el estilo texto theorem para cada tipo definido anteriormente
\newcounter{dummy} 
\numberwithin{dummy}{section}
\theoremstyle{ocrenumbox}
\newtheorem{theoremeT}[dummy]{Teorema}  % (Pedro: Theorem)
\newtheorem{problem}{Problema}[chapter]  % (Pedro: Problem)
\newtheorem{exerciseT}{Ejercicio}[chapter] % (Pedro: Exercise)
\theoremstyle{blacknumex}
\newtheorem{exampleT}{Ejemplo}[chapter] % (Pedro: Example)
\theoremstyle{blacknumbox}
\newtheorem{vocabulary}{Vocabulario}[chapter]  % (Pedro: Vocabulary)
\newtheorem{definitionT}{Definición}[section]  % (Pedro: Definition)
\newtheorem{corollaryT}[dummy]{Corolario}  % (Pedro: Corollary)
\theoremstyle{ocrenum}
\newtheorem{proposition}[dummy]{Proposición} % (Pedro: Proposition)


\usepackage[framemethod=default]{mdframed}



\newcommand{\intoo}[2]{\mathopen{]}#1\,;#2\mathclose{[}}
\newcommand{\ud}{\mathop{\mathrm{{}d}}\mathopen{}}
\newcommand{\intff}[2]{\mathopen{[}#1\,;#2\mathclose{]}}
\newtheorem{notation}{Notation}[chapter]


\mdfdefinestyle{exampledefault}{%
rightline=true,innerleftmargin=10,innerrightmargin=10,
frametitlerule=true,frametitlerulecolor=green,
frametitlebackgroundcolor=yellow,
frametitlerulewidth=2pt}


% Theorem box
\newmdenv[skipabove=7pt,
skipbelow=7pt,
backgroundcolor=black!5,
linecolor=ocre,
innerleftmargin=5pt,
innerrightmargin=5pt,
innertopmargin=10pt,%5pt
leftmargin=0cm,
rightmargin=0cm,
innerbottommargin=5pt]{tBox}

% Exercise box	  
\newmdenv[skipabove=7pt,
skipbelow=7pt,
rightline=false,
leftline=true,
topline=false,
bottomline=false,
backgroundcolor=ocre!10,
linecolor=ocre,
innerleftmargin=5pt,
innerrightmargin=5pt,
innertopmargin=10pt,%5pt
innerbottommargin=5pt,
leftmargin=0cm,
rightmargin=0cm,
linewidth=4pt]{eBox}	

% Definition box
\newmdenv[skipabove=7pt,
skipbelow=7pt,
rightline=false,
leftline=true,
topline=false,
bottomline=false,
linecolor=ocre,
innerleftmargin=5pt,
innerrightmargin=5pt,
innertopmargin=10pt,%0pt
leftmargin=0cm,
rightmargin=0cm,
linewidth=4pt,
innerbottommargin=0pt]{dBox}	

% Corollary box
\newmdenv[skipabove=7pt,
skipbelow=7pt,
rightline=false,
leftline=true,
topline=false,
bottomline=false,
linecolor=gray,
backgroundcolor=black!5,
innerleftmargin=5pt,
innerrightmargin=5pt,
innertopmargin=10pt,%5pt
leftmargin=0cm,
rightmargin=0cm,
linewidth=4pt,
innerbottommargin=5pt]{cBox}

% Crea un entorno para cada tipo de theorem y le asigna un estilo 
% con ayuda de las cajas coloreadas anteriores
\newenvironment{theorem}{\begin{tBox}\begin{theoremeT}}{\end{theoremeT}\end{tBox}}
\newenvironment{exercise}{\begin{eBox}\begin{exerciseT}}{\hfill{\color{ocre}\tiny\ensuremath{\blacksquare}}\end{exerciseT}\end{eBox}}				  
\newenvironment{definition}{\begin{dBox}\begin{definitionT}}{\end{definitionT}\end{dBox}}	
\newenvironment{example}{\begin{exampleT}}{\hfill{\tiny\ensuremath{\blacksquare}}\end{exampleT}}		
\newenvironment{corollary}{\begin{cBox}\begin{corollaryT}}{\end{corollaryT}\end{cBox}}	

%	ENVIRONMENT remark
\newenvironment{remark}{\par\vspace{10pt}\small 
% Espacio blanco vertical sobre la nota y tamaño de fuente menor
\begin{list}{}{
\leftmargin=35pt % Indentación sobre la izquierda
\rightmargin=25pt}\item\ignorespaces % Indentación sobre la derecha
\makebox[-2.5pt]{\begin{tikzpicture}[overlay]
\node[draw=ocre!60,line width=1pt,circle,fill=ocre!25,font=\sffamily\bfseries,inner sep=2pt,outer sep=0pt] at (-15pt,0pt){\textcolor{ocre}{N}}; \end{tikzpicture}} % R naranja en un círculo (Pedro)
\advance\baselineskip -1pt}{\end{list}\vskip5pt} 
% Espaciado de línea más estrecho y espacio en blanco después del comentario


\newenvironment{solutionExe}{\par\vspace{10pt}\small 
\begin{list}{}{
\leftmargin=35pt 
\rightmargin=25pt}\item\ignorespaces 
\makebox[-2.5pt]{\begin{tikzpicture}[overlay]
\node[draw=ocre!60,line width=1pt,circle,fill=ocre!25,font=\sffamily\bfseries,inner sep=2pt,outer sep=0pt] at (-15pt,0pt){\textcolor{ocre}{S}}; \end{tikzpicture}} 
\advance\baselineskip -1pt}{\end{list}\vskip5pt} 

\newenvironment{solutionExa}{\par\vspace{10pt}\small 
\begin{list}{}{
\leftmargin=35pt 
\rightmargin=25pt}\item\ignorespaces 
\makebox[-2.5pt]{\begin{tikzpicture}[overlay]
\node[draw=ocre!60,line width=1pt,circle,fill=ocre!55,font=\sffamily\bfseries,inner sep=2pt,outer sep=0pt] at (-15pt,0pt){\textcolor{ocre}{S}}; \end{tikzpicture}} 
\advance\baselineskip -1pt}{\end{list}\vskip5pt} 

\usepackage{tcolorbox}
\usepackage{multirow}

\usetikzlibrary{trees}

\theoremstyle{ocrenum}
\newtheorem{solutionT}[dummy]{Solución}  % (Pedro: Corollary)
\newenvironment{solution}{\begin{cBox}\begin{solutionT}}{\end{solutionT}\end{cBox}}	


\newcommand{\tcolorboxsolucion}[2]{%
\begin{tcolorbox}[colback=green!5!white,colframe=green!75!black,title=#1] 
 #2
 %\tcblower  % pone una línea discontinua
\end{tcolorbox}
}% final definición comando

\newtcbox{\mybox}[1][green]{on line,
arc=0pt,outer arc=0pt,colback=#1!10!white,colframe=#1!50!black, boxsep=0pt,left=1pt,right=1pt,top=2pt,bottom=2pt, boxrule=0pt,bottomrule=1pt,toprule=1pt}



\mdfdefinestyle{exampledefault}{%
rightline=true,innerleftmargin=10,innerrightmargin=10,
frametitlerule=true,frametitlerulecolor=green,
frametitlebackgroundcolor=yellow,
frametitlerulewidth=2pt}





\newcommand{\betheorem}{\begin{theorem}}
\newcommand{\eetheorem}{\end{theorem}}
\newcommand{\bedefinition}{\begin{definition}}
\newcommand{\eedefinition}{\end{definition}}

\newcommand{\beremark}{\begin{remark}}
\newcommand{\eeremark}{\end{remark}}
\newcommand{\beexercise}{\begin{exercise}}
\newcommand{\eeexercise}{\end{exercise}}
\newcommand{\beexample}{\begin{example}}
\newcommand{\eeexample}{\end{example}}
\newcommand{\becorollary}{\begin{corollary}}
\newcommand{\eecorollary}{\end{corollary}}


\newcommand{\besolutionExe}{\begin{solutionExe}}
\newcommand{\eesolutionExe}{\end{solutionExe}}
\newcommand{\besolutionExa}{\begin{solutionExa}}
\newcommand{\eesolutionExa}{\end{solutionExa}}


%%%%%%%%


% Caja Salida Markdown
\newmdenv[skipabove=7pt,
skipbelow=7pt,
rightline=false,
leftline=true,
topline=false,
bottomline=false,
backgroundcolor=GreenYellow!10,
linecolor=GreenYellow!80,
innerleftmargin=5pt,
innerrightmargin=5pt,
innertopmargin=10pt,%5pt
innerbottommargin=5pt,
leftmargin=0cm,
rightmargin=0cm,
linewidth=4pt]{mBox}	

%% RMarkdown
\newenvironment{markdownsal}{\begin{mBox}}{\end{mBox}}	

\newcommand{\bmarkdownsal}{\begin{markdownsal}}
\newcommand{\emarkdownsal}{\end{markdownsal}}


\usepackage{array}
\usepackage{multirow}
\usepackage{wrapfig}
\usepackage{colortbl}
\usepackage{pdflscape}
\usepackage{tabu}
\usepackage{threeparttable}
\usepackage{subfig} %new
%\usepackage{booktabs,dcolumn,rotating,thumbpdf,longtable}
\usepackage{dcolumn,rotating}  %new
\usepackage[graphicx]{realboxes} %new de: https://stackoverflow.com/questions/51633434/prevent-pagebreak-in-kableextra-landscape-table

%define el interlineado vertical
%\renewcommand{\baselinestretch}{1.5}

%define etiqueta para las Tablas o Cuadros
%\renewcommand\spanishtablename{Tabla}

%%\bibliographystyle{plain} %new no necesario


%%%%%%%%%%%% PARA USO CON biblatex
% \DefineBibliographyStrings{english}{%
%   backrefpage = {ver pag.\adddot},%
%   backrefpages = {ver pags.\adddot}%
% }

% \DefineBibliographyStrings{spanish}{%
%   backrefpage = {ver pag.\adddot},%
%   backrefpages = {ver pags.\adddot}%
% }
% 
% \DeclareFieldFormat{pagerefformat}{\mkbibparens{{\color{red}\mkbibemph{#1}}}}
% \renewbibmacro*{pageref}{%
%   \iflistundef{pageref}
%     {}
%     {\printtext[pagerefformat]{%
%        \ifnumgreater{\value{pageref}}{1}
%          {\bibstring{backrefpages}\ppspace}
%          {\bibstring{backrefpage}\ppspace}%
%        \printlist[pageref][-\value{listtotal}]{pageref}}}}
% 
%%% de kableExtra
\usepackage{booktabs}
\usepackage{longtable}
%\usepackage{array}
%\usepackage{multirow}
%\usepackage{wrapfig}
%\usepackage{float}
%\usepackage{colortbl}
%\usepackage{pdflscape}
%\usepackage{tabu}
%\usepackage{threeparttable}
\usepackage{threeparttablex}
\usepackage[normalem]{ulem}
\usepackage{makecell}
%\usepackage{xcolor}

%%%%%%% fin: latex_preambulo.tex PLLC


\begin{document}

\bibliographystyle{flexbib}



\raggedbottom

\ifdefined\ifprincipal
\else
\setlength{\parindent}{1em}
\pagestyle{fancy}
\setcounter{tocdepth}{4}
\tableofcontents

\fi

\ifdefined\ifdoblecara
\fancyhead{}{}
\fancyhead[LE,RO]{\scriptsize\rightmark}
\fancyfoot[LO,RE]{\scriptsize\slshape \leftmark}
\fancyfoot[C]{}
\fancyfoot[LE,RO]{\footnotesize\thepage}
\else
\fancyhead{}{}
\fancyhead[RO]{\scriptsize\rightmark}
\fancyfoot[LO]{\scriptsize\slshape \leftmark}
\fancyfoot[C]{}
\fancyfoot[RO]{\footnotesize\thepage}
\fi

\renewcommand{\headrulewidth}{0.4pt}
\renewcommand{\footrulewidth}{0.4pt}

\chapter{Concepto de solución}\label{concepto-de-soluciuxf3n}

En la teoría de juegos, un \textbf{concepto de solución} es una regla
formal para predecir cómo se jugará un juego. Estas predicciones se
denominan ``soluciones'', y describen las estrategias que adoptarán los
jugadores y, por tanto, el resultado que puede esperarse del juego.

En general, los conceptos de solución tratan de identificar perfiles de
estrategias que puedan considerarse racionales o estables, teniendo en
cuenta que cada jugador toma sus decisiones considerando las posibles
acciones de los demás participantes.

Los conceptos de solución más comúnmente utilizados son conceptos de
equilibrio, el más famoso de ellos es el equilibrio de Nash.

En muchos juegos, un mismo concepto de solución puede dar lugar a varias
soluciones posibles. Por este motivo, la teoría de juegos introduce en
ocasiones refinamientos de equilibrio, cuyo objetivo es reducir el
conjunto de soluciones y seleccionar aquellas que resultan más
razonables.

A continuación, se presentan algunos conceptos de solución:

\section{Estrategias dominadas y
dominantes}\label{estrategias-dominadas-y-dominantes}

Una estrategia de un jugador se dice dominante si es tan buena o más que
cualquier otra como respuesta a cualquier combinación de estrategias que
elijan los demás jugadores, y una estrategia dada \(s_i\) de un jugador
se dice que está dominada por otra estrategia \(s'_i\) del mismo jugador
si la segunda le conviene más que la primera, independientemente de lo
que hagan los otros jugadores. El argumento básico de dominación
consiste en que un jugador racional no debería jugar estrategias
dominadas y en que, en caso de saber que otros jugadores son racionales,
debería suponer que éstos no van a jugar tal clase de estrategias.

\textbf{Definición 3.1.} En el juego
\(G = \{S_1, \ldots, S_n; u_1,\ldots,u_n\}\), sean \(s'_i\) y \(s''_i\)
dos estrategias del jugador \emph{i}.

\begin{enumerate}
\def\labelenumi{\alph{enumi})}
\tightlist
\item
  Decimos que \(s'_i\) está \textbf{dominada}, o también
  \textbf{débilmente dominada}, por \(s''_i\) cuando la desigualdad
  \[u_i(s_1,\ldots,s_{i-1},s'_i,s_{i+1},\ldots,s_n)\le u_i(s_1,\ldots,s_{i-1},s''_i,s_{i+1},\ldots,s_n)\]
\end{enumerate}

se cumple para toda combinación de estrategias \(s_{-i}\) de los otros
jugadores, y para alguna de esas combinaciones se cumple de modo
estricto. Decimos de manera equivalente que \(s''_i\) \textbf{domina} a
\(s'_i\). (Es decir, siempre le conviene usar \(s''_i\) al menos tanto
como usar \(s'_i\), hagan lo que hagan los otros jugadores, y a veces le
conviene más)

\begin{enumerate}
\def\labelenumi{\alph{enumi})}
\setcounter{enumi}{1}
\item
  Decimos que \(s'_i\) está \textbf{estrictamente dominada} por
  \(s''_i\) cuando la desigualdad
  \[u_i(s_1,\ldots,s_{i-1},s'_i,s_{i+1},\ldots,s_n)<u_i(s_1,\ldots,s_{i-1},s''_i,s_{i+1},\ldots,s_n)\]
  se cumple para toda combinación de estrategias \(s_{-i}\) de los otros
  jugadores. (Es decir, le conviene más usar \(s''_i\) que \(s'_i\),
  hagan lo que hagan los otros jugadores)
\item
  Decimos que \(s'_i\) es \textbf{no dominada} si no existe ninguna otra
  estrategia del jugador \emph{i} que la domine, y decimos que \(s'_i\)
  es \textbf{no dominada estrictamente} si no existe ninguna otra
  estrategia del jugador i que la domine estrictamente.
\end{enumerate}

Por otra parte, el concepto de estrategia dominante que se define a
continuación, se aplica a una estrategia de un jugador cuando ésta es
tan buena o más que cualquier otra como respuesta a cualquier
combinación de estrategias qie elijan los demás jugadores.

\textbf{Definición 3.2.} En el juegos
\(G = \{S_1,\ldots,S_n;u_1,\ldots,u_n\}\), sea \(s'_i\) una estrategia
del jugador \emph{i}.

\begin{enumerate}
\def\labelenumi{\alph{enumi})}
\item
  Decimos qie \(s'_i\) es \textbf{dominante} cuando la desigualdad
  \[u_i(s_1,\ldots,s_{i-1},s_i,s_{i+1},\ldots,s_n)\le u_i(s_1,\ldots,s_{i-1},s'_i,s_{i+1},\ldots,s_n)\]
  se cumple para toda estrategia \(s_i\) de dicho jugador y para toda
  combinación de estrategias \(s_{-i}\) de los otros jugadores. (Es
  decir, siempre le conviene usar la estrategia \(s'_i\) al menos tanto
  como cualquier otra, hagan lo que hagan los otros jugadores)
\item
  Si todas las desigualdades se cumplen de manera estricta (para
  \(s_i \neq s'_i\)), decimos que \(s'_i\) es \textbf{estrictamente
  dominante}.
\end{enumerate}

\begin{tcolorbox}[colback=blue!10]
\textbf{Ejemplo 3.1} \\

Consideremos el siguiente juego de dos jugadores (Dilema del Prisionero)

\begin{table}[H]
\centering
\begin{tabular}{c|cc}
 & \multicolumn{2}{c}{\textbf{Preso 2}} \\ 
\textbf{Preso 1} & Callar & Confesar \\ \hline
Callar & (4,4) & (0,5) \\
Confesar & (5,0) & (1,1)
\end{tabular}
\caption{Pagos del Dilema del Prisionero}
\label{tab:dilema_prisionero}
\end{table}

\textbf{Análisis:}  

Para el jugador 1:

\[
u_1(\text{Callar, Callar}) = 4 < 5 = u_1(\text{Confesar, Callar})
\]

\[
u_1(\text{Callar, Confesar}) = 0 < 1 = u_1(\text{Confesar, Confesar})
\]

De manera análoga, para el jugador 2:

\[
u_2(\text{Callar, Callar}) = 4 < 5 = u_2(\text{Callar, Confesar})
\]

\[
u_2(\text{Confesar, Callar}) = 0 < 1 = u_2(\text{Confesar, Confesar})
\]

\textbf{Conclusión:}  

Para ambos jugadores, la estrategia \textbf{Callar} está \textbf{estrictamente dominada} por la estrategia \textbf{Confesar}.  
Por lo tanto, si los jugadores son racionales, ambos elegirán \textbf{Confesar} y el equilibrio estrictamente dominante es \((\text{Confesar}, \text{Confesar})\).
\end{tcolorbox}

\subsection{Uso de Estrategias Dominantes
(UED)}\label{uso-de-estrategias-dominantes-ued}

Según este concepto de solución, pertenecen a la solución del juego
todos aquellos perfiles de estrategias en los cuales cada jugador usa
una estrategia dominante.

Naturalmente, en caso de ser aplicable, este concepto de solución es
obvio, pues cualquier jugador racional jugará una estrategia dominante
si dispone de ella, independientemente de cualquier otra consideración
(y si tiene varias, jugará una de ellas). En particular lo hará
independientemente de lo que sepa o suponga acerca de cómo van a jugar
los otros jugadores, y de qué sepa acerca de las características de
dichos jugadores. De hecho, incluso cabe la posibilidad de que los
jugadores tengan un conocimiento muy limitado de aquellos aspectos del
juego que afectan a los demás jugadores (por ejemplo, podrían no conocer
los pagos de los otros), y aun así tender a jugar su estrategia
dominante.

Por desgracia, este concepto de solución no siempre es aplicable. Los
juegos en los que cada jugador tiene alguna estrategia dominante son más
bien la excepción que la regla.

\begin{tcolorbox}[colback=blue!10]
\textbf{Ejemplo 3.2} \\

En el dilema del prisionero, este concepto sí es aplicable, pues en este caso cada jugador tiene una estrategia estrictamente dominante, que es confesar.

Mientras en el siguiente ejemplo  

\begin{table}[H]
\centering
\begin{tabular}{c|ccc}
 & \multicolumn{3}{c}{\textbf{Jugador 2}} \\ 
\textbf{Jugador 1} & I & C & D \\ \hline
A & (3,y) & (4,2) & (1,x) \\
M & (2,4) & (3,5) & (4,0) \\
B & (1,0) & (2,1) & (0,3)
\end{tabular}
\caption{Juego con estrategias estrictamente dominadas pero sin estrategias dominantes}
\label{tab:contraejemplo_UED}
\end{table}

\textbf{Conclusión:}  
En este juego, con x = y = 1, la estrategia B del jugador 1 (J1) está estrictamente dominada por las estrategias A y M. Pero ni A ni M son estrategias dominantes. Por otra parte, para el jugador 2 (J2) la estrategia I está estrictamente dominada por la estrategia C, pero ni C ni D son dominantes. Así pues, aunque el concepto de solución anterior no es aplicable, el argumento básico de dominación (ningún jugador racional debería jugar estrategias dominadas) nos permite concluir inmediatamente que ni J1 jugará B ni J2 jugará I, pero sólo nos permite esa conclusión
\end{tcolorbox}

\subsection{Eliminación Iterativa Estricta
(EIE)}\label{eliminaciuxf3n-iterativa-estricta-eie}

\textbf{Definición 3.3.} Dado un juego finito o infinito
\(G = \{S_1, \ldots, S_n; u_1, \ldots, u_n\}\), llamamos
\textbf{Eliminación Iterativa Estricta (EIE)}, o bien
\textbf{Eliminación Iterativa de Estrategias Estrictamente Dominadas},
al siguiente proceso de eliminación:

\textbf{1º Paso.} De cada uno de los jugadores, y a la vez, se eliminan
todas las estrategias que estén estrictamente dominadas en el juego
inicial G. Se construye el juego reducido \(G_1\) que resulta de tal
eliminación.

\textbf{2º Paso.} De cada uno de los jugadores, y a la vez, se eliminan
todas las estrategias que estén estrictamente dominadas en el juego
reducido \(G_1\). Se construye el juego reducido \(G_2\) que resulta de
tal eliminación.

Y así sucesivamente. Se acaba el proceso cuando ya no quedan, para
ningún jugador, estrategias que eliminar. Denotamos \(S^S_i\) al
conjunto de las estrategias supervivientes del jugador i, y las llamamos
estrategias iterativamente no dominadas.

De acuerdo con este concepto de solución, son soluciones todos los
perfiles estratégicos constituidos por estrategias que sobreviven al
proceso de eliminación iterativa estricta. LLamaremos
\(\mathbf{S^{EIE}}\) al conjunto de dichos perfiles estratégicos
solución.

\begin{tcolorbox}[colback=blue!10]
\textbf{Ejemplo 3.3} \\

Dado el juego

\begin{table}[H]
\centering
\begin{tabular}{c|cc}
 & \multicolumn{2}{c}{\textbf{Jugador 2}} \\ 
\textbf{Jugador 1} & Izquierda & Derecha\\ \hline
Alta & (4,2) & (0,1) \\
Media & (1,2) & (2,4) \\
Baja & (3,3) & (4,2)
\end{tabular}
\caption{Juego G (EIE)}
\label{tab:EIE}
\end{table}

Puesto que la racionalidad de ambos jugadores es una información de dominio público, ambos jugadores son racionales y cada uno sabe que el otro lo es y además sabe que el otro lo sabe. En consecuencia, puede razonarse así:

El jugador 1 (J1), que es racional, elimina la estrategia \textbf{Media} porque está estrictamente dominada por \textbf{Baja} (pagos 1 frente a 3 y 2 frente a 4). El jugador 2 (J2) sabe que J1 es racional y por tanto sabe que J1 ha eliminado \textbf{Media}. Al comparar \textbf{Izquierda} con \textbf{Derecha} en ausencia de Media de J1, el jugador J2, que también es racional, deduce que \textbf{Izquierda} domina estrictamente a \textbf{Derecha} (pagos 2 frente a 1 y 3 frente a 2), por lo que elimina \textbf{Derecha}.

\begin{table}[H]
\centering
\begin{tabular}{c|c}
 & \multicolumn{1}{c}{\textbf{Jugador 2}} \\ 
\textbf{Jugador 1} & Izquierda\\ \hline
Alta & (4,2)\\
Baja & (3,3)
\end{tabular}
\caption{Juego $G_1$ (EIE)}
\label{tab:EIE}
\end{table}

El jugador J1 sabe que J2 es racional y que además J2 sabe que J1 es racional, y por tanto sabe que J2 ha eliminado \textbf{Derecha} como consecuencia de la eliminación por J1 de \textbf{Media}. Al comparar \textbf{Alta} con \textbf{Baja} en ausencia de \textbf{Derecha}, J1 deduce que \textbf{Alta} domina estrictamente a \textbf{Baja} (pago 4 frente a 3) por lo que elimina \textbf{Baja}.

\begin{table}[H]
\centering
\begin{tabular}{c|c}
 & \multicolumn{1}{c}{\textbf{Jugador 2}} \\ 
\textbf{Jugador 1} & Izquierda\\ \hline
Alta & (4,2)\\
\end{tabular}
\caption{Juego $G_2$ perfil superviviente único (EIE)}
\label{tab:EIE}
\end{table}

En conclusión, las únicas estrategias supervivientes al proceso EIE son \textbf{Alta} de J1 e \textbf{Izquierda} de J2. Ambas constituyen el único perfil solución del juego, es decir, $\mathbf{S^{EIE}}$ = \{(Alta, Izquierda)\}
\end{tcolorbox}

\subsection{Eliminación Iterativa Débil
(EID)}\label{eliminaciuxf3n-iterativa-duxe9bil-eid}

\textbf{Definición 3.4.} Dado un juego finito o infinito
\(G = \{S_1,\ldots,S_n;u_1,\ldots,u_n\}\), llamamos \textbf{Eliminación
Iterativa Débil (EID)}, o bien \textbf{Eliminación Iterativa de
Estrategias Débilmente Dominadas}, al proceso de eliminación siguiente:

\textbf{1º Paso.} De cada uno de los jugadores, y a la vez, se elimina
todas las estrategias que estén débilmente dominadas en el juego inicial
G. Se construye el juego reducido \(G_1\) que resulta de tal
eliminación.

\textbf{2º Paso.} De cada uno de los jugadores, y a la vez, se eliminan
todas las estrategias que estén débilmente dominadas en el juego
reducido \(G_1\). Se construye el juego reducido \(G_2\) que resulta de
tal eliminación.

Y así sucesivamente. Se acaba el proceso cuando ya no quedan, para
ningún jugador i, estrategias que eliminar. Denotamos \(S^S_i\) al
conjunto de las estrategias supervivientes del jugador i, y las llamamos
estrategias iterativamente no dominadas.

Desgraciadamente, y al contrario de lo que ocurre para los algoritmos de
eliminación estricta, ahora el resultado final del proceso de
eliminación sí depende del algoritmo utilizado, es decir, del orden y
modo en que se han ido seleccionando las estrategias a eliminar.

Este concepto de solución determina como solución \(\mathbf{S^{EID}}\).

\begin{tcolorbox}[colback=blue!10]
\textbf{Ejemplo 3.4} \\

Dado el juego

\begin{table}[H]
\centering
\begin{tabular}{c|ccc}
 & \multicolumn{3}{c}{\textbf{Jugador 2}} \\ 
\textbf{Jugador 1} & I & C & D\\ \hline
A & (3,1) & (4,2) & (1,2)\\
M & (2,4) & (3,5) & (4,0) \\
B & (1,0) & (2,1) & (0,3)
\end{tabular}
\caption{Juego G (EID)}
\label{tab:EID}
\end{table}

Puesto que la racionalidad de ambos jugadores es información de dominio público, ambos jugadores son racionales y cada uno sabe que el otro lo es y además sabe que el otro lo sabe. En consecuencia, puede razonarse del siguiente modo:

El jugador 1 (J1), que es racional, compara sus estrategias \textbf{A}, \textbf{M} y \textbf{B}. Al comparar \textbf{A} con \textbf{B}, observa que \textbf{A} proporciona siempre un pago mayor que \textbf{B} independientemente de la estrategia elegida por el jugador 2 (pagos 3 frente a 1 si J2 juega \textbf{I}, 4 frente a 2 si J2 juega \textbf{C}, y 1 frente a 0 si J2 juega \textbf{D}). Por tanto, la estrategia \textbf{B} está estrictamente dominada por \textbf{A} y J1 elimina \textbf{B}.

El jugador 2 (J2) sabe que J1 es racional y por tanto sabe que J1 ha eliminado la estrategia \textbf{B}. En consecuencia, J2 compara ahora sus estrategias considerando únicamente las estrategias \textbf{A} y \textbf{M} de J1. Al comparar \textbf{I} con \textbf{C}, observa que \textbf{C} proporciona siempre un pago mayor que \textbf{I} (pagos 2 frente a 1 si J1 juega \textbf{A}, y 5 frente a 4 si J1 juega \textbf{M}). Por tanto, \textbf{C} domina estrictamente a \textbf{I} y J2 elimina la estrategia \textbf{I}.

\end{tcolorbox}

\begin{tcolorbox}[colback=blue!10]
\begin{table}[H]
\centering
\begin{tabular}{c|cc}
 & \multicolumn{2}{c}{\textbf{Jugador 2}} \\ 
\textbf{Jugador 1} & C & D\\ \hline
A & (4,2) & (1,2)\\
M & (3,5) & (4,0)
\end{tabular}
\caption{Juego $G_1$ (EID)}
\label{tab:EID}
\end{table}

El jugador J1 sabe que J2 es racional y que además J2 sabe que J1 es racional, por lo que sabe que J2 ha eliminado \textbf{I}. Considerando el juego reducido, vemos que para J1 ninguna estrategia domina a otra, es decir, no hay dominancia débil, mientras que J2 compara ahora las estrategias \textbf{C} y \textbf{D}. Observa que \textbf{C} proporciona un pago al menos igual que \textbf{D} cuando J1 juega \textbf{A} (pagos 2 frente a 2) y un pago mayor cuando J1 juega \textbf{M} (pagos 5 frente a 0). Por tanto, \textbf{D} está débilmente dominada por \textbf{C}, y J2 elimina la estrategia \textbf{D}.

\begin{table}[H]
\centering
\begin{tabular}{c|c}
 & \multicolumn{1}{c}{\textbf{Jugador 2}} \\ 
\textbf{Jugador 1} & C\\ \hline
A & (4,2)\\
M & (3,5)
\end{tabular}
\caption{Juego $G_2$ (EID)}
\label{tab:EID}
\end{table}

Finalmente, el jugador J1 compara sus estrategias \textbf{A} y \textbf{M}. Observa que \textbf{A} proporciona un pago mayor que \textbf{M} (4 frente a 3), por lo que \textbf{M} está dominada por \textbf{A}. En consecuencia, J1 elimina la estrategia \textbf{M}.

\begin{table}[H]
\centering
\begin{tabular}{c|c}
 & \multicolumn{1}{c}{\textbf{Jugador 2}} \\ 
\textbf{Jugador 1} & C\\ \hline
A & (4,2)
\end{tabular}
\caption{Juego $G_3$ perfil superviviente único (EID)}
\label{tab:EID}
\end{table}

En conclusión, las únicas estrategias supervivientes al proceso de eliminación iterativa de estrategias débilmente dominadas son \textbf{A} para el jugador 1 y \textbf{C} para el jugador 2. Ambas constituyen el único perfil solución del juego, es decir, $\mathbf{S^{EID}}$ = \{(A,C)\}.
\end{tcolorbox}

Cabe mencionar que todos estos métods son evidentemente aplicables a
juegos expresados en su forma extensiva:

\begin{center}
\begin{tikzpicture}[
scale=1.2,
level distance=2.8cm,
level 1/.style={sibling distance=7cm},
level 2/.style={sibling distance=2.8cm}
]

\node[circle,draw]{J1}
child {node[circle,draw]{J2}
    child {node[rectangle,draw]{(3,1)} edge from parent node[left]{I}}
    child {node[rectangle,draw]{(4,2)} edge from parent node[above]{C}}
    child {node[rectangle,draw]{(1,2)} edge from parent node[right]{D}}
    edge from parent node[above left]{A}
}
child {node[circle,draw]{J2}
    child {node[rectangle,draw]{(2,4)} edge from parent node[left]{I}}
    child {node[rectangle,draw]{(3,5)} edge from parent node[above]{C}}
    child {node[rectangle,draw]{(4,0)} edge from parent node[right]{D}}
    edge from parent node[above]{M}
}
child {node[circle,draw]{J2}
    child {node[rectangle,draw]{(1,0)} edge from parent node[left]{I}}
    child {node[rectangle,draw]{(2,1)} edge from parent node[above]{C}}
    child {node[rectangle,draw]{(0,3)} edge from parent node[right]{D}}
    edge from parent node[above right]{B}
};

\end{tikzpicture}
\end{center}

\subsection{Juegos resolubles por
dominación}\label{juegos-resolubles-por-dominaciuxf3n}

\section{Equilibrio de Nash}\label{equilibrio-de-nash}

\subsection{Equilibrio de Nash en estrategias
puras}\label{equilibrio-de-nash-en-estrategias-puras}

\subsection{Equilibrio de Nash en estrategias
mixtas}\label{equilibrio-de-nash-en-estrategias-mixtas}

\section{Optimalidad de Pareto}\label{optimalidad-de-pareto}

\section{Equilibrio perfecto en
subjuegos}\label{equilibrio-perfecto-en-subjuegos}

\bibliography{bib/library.bib,bib/paquetes.bib}


%


\end{document}
